%% AUTO-GENERATED by hrdaya.latex_gen
%% Data version: 1.2.0  Hash: b0f93d2af26b
%% Generated: 2026-02-11T23:12:22Z
%% Regenerate with: make latex-gen

\documentclass[10pt,landscape,a4paper]{article}

\usepackage[landscape,margin=1.2cm,bottom=2cm]{geometry}
\usepackage{fontspec}
\usepackage{xeCJK}
\usepackage{tabularx}
\usepackage{array}
\usepackage{booktabs}
\usepackage{xcolor}
\usepackage{longtable}
\usepackage{multicol}
\usepackage[perpage]{footmisc}
\usepackage{fancyhdr}
\pagestyle{fancy}
\fancyhf{}
\renewcommand{\headrulewidth}{0.4pt}
\fancyhead[L]{\small Prajñāpāramitāhṛdaya -- Combined Critical Edition}
\fancyhead[R]{\small 般若波羅蜜多心經}
\fancyfoot[C]{\thepage}
\usepackage{hyperref}
\hypersetup{colorlinks=true,linkcolor=blue}

\setmainfont{Noto Serif}[
  BoldFont = {Noto Serif},
  BoldFeatures = {RawFeature=+axis{wght=700}},
]
\setsansfont{Noto Sans}
\setCJKmainfont{Noto Serif CJK SC}
\setCJKsansfont{Noto Sans CJK SC}

\newfontfamily\devanagarifont{Noto Serif Devanagari}[
  BoldFont = {Noto Serif Devanagari},
  BoldFeatures = {RawFeature=+axis{wght=700}},
]
\newcommand{\deva}[1]{{\devanagarifont #1}}

\newfontfamily\tibetanfont{Noto Serif Tibetan}[
  BoldFont = {Noto Serif Tibetan},
  BoldFeatures = {RawFeature=+axis{wght=700}},
]
\newcommand{\tib}[1]{{\tibetanfont\XeTeXlinebreaklocale "bo" #1}}

\newcommand{\segnum}[1]{\textbf{[#1]}}
\newcommand{\rom}[1]{{\small\color{gray} #1}}
\newcommand{\engltrans}[1]{\multicolumn{3}{p{0.95\textwidth}}{\small\itshape #1} \\}
\newcommand{\var}[1]{{\color{red}\textsuperscript{\textbf{#1}}}}
\newcommand{\lemma}[1]{\textbf{#1}}
\newcommand{\reading}[1]{#1}
\newcommand{\wit}[1]{\textsc{#1}}
\newcommand{\depmark}[2]{{\scriptsize\color{blue}[#1$\to$#2]}}

\begin{document}
\begin{center}
{\Huge\bfseries Prajñāpāramitāhṛdaya}\\[0.3em]
{\Huge 般若波羅蜜多心經}\\[0.3em]
{\huge\deva{प्रज्ञापारमिताहृदय}}\\[0.5em]
{\LARGE Heart Sūtra}\\[1em]
{\Large Combined Critical Edition}\\[0.5em]
{\normalsize Chinese · Sanskrit · Tibetan}\\[2em]
\end{center}

\hrule
\vspace{1em}


\footnotesize

\begin{longtable}{|>{\raggedright\arraybackslash}p{0.30\textwidth}|>{\raggedright\arraybackslash}p{0.36\textwidth}|>{\raggedright\arraybackslash}p{0.28\textwidth}|}
\hline
\textbf{Chinese (T251)} & \textbf{Sanskrit} & \textbf{Tibetan (Toh 21)} \\
\hline
\hline
\endfirsthead

\hline
\textbf{Chinese (T251)} & \textbf{Sanskrit} & \textbf{Tibetan (Toh 21)} \\
\hline
\hline
\endhead

\hline
\multicolumn{3}{r}{\small\itshape continued...}
\endfoot

\hline
\endlastfoot

\segnum{1} 觀自在菩薩，行深般若波羅蜜多時，照見五蘊皆空，度一切苦厄。
\newline\rom{Guānzìzài púsà, xíng shēn bōrě bōluómìduō shí, zhàojiàn wǔyùn jiē kōng, dù yīqiè kǔ è.}
&
\deva{आर्य-अवलोकितेश्वरो बोधिसत्त्वो गम्भीरां प्रज्ञापारमिताचर्यां चरमाणो व्यवलोकयति स्म: पञ्च-स्कन्धास्, तांश्च स्वभाव-शून्यान् पश्यति स्म।}
\newline\rom{Ārya-avalokiteśvaro bodhisattvo gambhīrāṃ prajñāpāramitācaryāṃ caramāṇo vyavalokayati sma: pañca-skandhās, tāṃś ca svabhāva-śūnyān paśyati sma.}
&
\tib{ཡང་དེའི་ཚེ་བྱང་ཆུབ་སེམས་དཔའ་སེམས་དཔའ་ཆེན་པོ་འཕགས་པ་སྤྱན་རས་གཟིགས་དབང་ཕྱུག་ཤེས་རབ་ཀྱི་ཕ་རོལ་ཏུ་ཕྱིན་པ་ཟབ་མོའི་སྤྱོད་པ་ཉིད་ལ་རྣམ་པར་བལྟ་ཞིང་། ཕུང་པོ་ལྔ་པོ་དེ་དག་ལ་ཡང་རང་བཞིན་གྱིས་སྟོང་པར་རྣམ་པར་བལྟའོ།}
\newline\rom{yang de'i tshe byang chub sems dpa' sems dpa' chen po 'phags pa spyan ras gzigs dbang phyug shes rab kyi pha rol tu phyin pa zab mo'i spyod pa nyid la rnam par blta zhing / phung po lnga po de dag la yang rang bzhin gyis stong par rnam par blta'o/}
\\
\hline
\engltrans{When Avalokiteśvara Bodhisattva was practicing the profound prajñāpāramitā, he illuminated the five aggregates and saw that they are all empty, thus transcending all suffering.}
\hline
\hline

\segnum{2} 舍利子，色不異空，空不異色；色即是空，空即是色。受想行識亦復如是。
\newline\rom{Shèlìzǐ, sè bù yì kōng, kōng bù yì sè; sè jí shì kōng, kōng jí shì sè. Shòu xiǎng xíng shí yì fù rúshì.}
&
\deva{इह शारिपुत्र रूपं शून्यता शून्यतैव रूपम्, रूपान् न पृथक् शून्यता शून्यतायाः न पृथग् रूपम्, यद् रूपं सा शून्यता या शून्यता तद् रूपम्; एवमेव वेदना-संज्ञा-संस्कार-विज्ञानानि।}
\newline\rom{Iha Śāriputra rūpaṃ śūnyatā śūnyataiva rūpam, rūpān na pṛthak śūnyatā śūnyatāyā na pṛthag rūpam, yad rūpaṃ sā śūnyatā yā śūnyatā tad rūpam; evam eva vedanā-saṃjñā-saṃskāra-vijñānāni.}
&
\tib{གཟུགས་སྟོང་པའོ། སྟོང་པ་ཉིད་གཟུགས་སོ། གཟུགས་ལས་སྟོང་པ་ཉིད་གཞན་མ་ཡིན། སྟོང་པ་ཉིད་ལས་ཀྱང་གཟུགས་གཞན་མ་ཡིན་ནོ། དེ་བཞིན་དུ་ཚོར་བ་དང་། འདུ་ཤེས་དང་། འདུ་བྱེད་དང་། རྣམ་པར་ཤེས་པ་རྣམས་སྟོང་པའོ།}
\newline\rom{gzugs stong pa'o/ stong pa nyid gzugs so/ gzugs las stong pa nyid gzhan ma yin/ stong pa nyid las kyang gzugs gzhan ma yin no/ de bzhin du tshor ba dang / 'du shes dang / 'du byed dang / rnam par shes pa rnams stong pa'o/}
\\
\hline
\engltrans{Śāriputra, form is not different from emptiness, emptiness is not different from form; form is emptiness, emptiness is form. Feeling, perception, formation, and consciousness are also like this.}
\hline
\hline

\segnum{3} 舍利子，是諸法空相，不生不滅，不垢不淨，不增不減。
\newline\rom{Shèlìzǐ, shì zhūfǎ kōng xiāng, bù shēng bù miè, bù gòu bù jìng, bù zēng bù jiǎn.}
&
\deva{इह शारिपुत्र सर्व-धर्माः शून्यता-लक्षणा, अनुत्पन्ना अनिरुद्धा, अमला अविमला, अनूना अपरिपूर्णाः।}
\newline\rom{Iha Śāriputra sarva-dharmāḥ śūnyatā-lakṣaṇā, anutpannā aniruddhā, amalā avimalā, anūnā aparipūrṇāḥ.}
&
\tib{ཤཱ་རིའི་བུ། དེ་ལྟར་ཆོས་ཐམས་ཅད་སྟོང་པ་ཉིད་དེ། མཚན་ཉིད་མེད་པ། མ་སྐྱེས་པ། མ་འགགས་པ། དྲི་མ་མེད་པ། དྲི་མ་དང་བྲལ་བ་མེད་པ། བྲི་བ་མེད་པ། གང་བ་མེད་པའོ།}
\newline\rom{shA ri'i bu/ de ltar chos thams cad stong pa nyid de/ mtshan nyid med pa/ ma skyes pa/ ma 'gags pa/ dri ma med pa/ dri ma dang bral ba med pa/ bri ba med pa/ gang ba med pa'o/}
\\
\hline
\engltrans{Śāriputra, the emptiness characteristic of all dharmas is: not arising, not ceasing; not defiled, not pure; not increasing, not decreasing.}
\hline
\hline

\segnum{4} 是故空中，無色，無受想行識，
\newline\rom{Shìgù kōng zhōng, wú sè, wú shòu xiǎng xíng shí,}
&
\deva{तस्माच् छारिपुत्र शून्यतायां न रूपं न वेदना न संज्ञा न संस्काराः न विज्ञानम्।}
\newline\rom{Tasmāc Chāriputra śūnyatāyāṃ na rūpaṃ na vedanā na saṃjñā na saṃskārāḥ na vijñānam.}
&
\tib{ཤཱ་རིའི་བུ། དེ་ལྟ་བས་ན་སྟོང་པ་ཉིད་ལ་གཟུགས་མེད། ཚོར་བ་མེད། འདུ་ཤེས་མེད། འདུ་བྱེད་རྣམས་མེད། རྣམ་པར་ཤེས་པ་མེད།}
\newline\rom{shA ri'i bu/ de lta bas na stong pa nyid la gzugs med/ tshor ba med/ 'du shes med/ 'du byed rnams med/ rnam par shes pa med/}
\\
\hline
\engltrans{Therefore, in emptiness there is no form, no feeling, perception, formation, or consciousness;}
\hline
\hline

\segnum{5} 無眼耳鼻舌身意，無色聲香味觸法，
\newline\rom{wú yǎn ěr bí shé shēn yì, wú sè shēng xiāng wèi chù fǎ,}
&
\deva{न चक्षुः-श्रोत्र-घ्राण-जिह्वा-काय-मनांसि। न रूप-शब्द-गन्ध-रस-स्प्रष्टव्य-धर्माः।}
\newline\rom{Na cakṣuḥ-śrotra-ghrāṇa-jihvā-kāya-manāṃsi. Na rūpa-śabda-gandha-rasa-spraṣṭavya-dharmāḥ.}
&
\tib{མིག་མེད། རྣ་བ་མེད། སྣ་མེད། ལྕེ་མེད། ལུས་མེད། ཡིད་མེད། གཟུགས་མེད། སྒྲ་མེད། དྲི་མེད། རོ་མེད། རེག་བྱ་མེད། ཆོས་མེད་དོ།}
\newline\rom{mig med/ rna ba med/ sna med/ lce med/ lus med/ yid med/ gzugs med/ sgra med/ dri med/ ro med/ reg bya med/ chos med do/}
\\
\hline
\engltrans{no eyes, ears, nose, tongue, body, or mind; no forms, sounds, smells, tastes, tangibles, or mental objects;}
\hline
\hline

\segnum{6} 無眼界，乃至無意識界。
\newline\rom{wú yǎn jiè, nǎizhì wú yìshí jiè.}
&
\deva{न चक्षुर्-धातुर् यावन् न मनोविज्ञान-धातुः।}
\newline\rom{Na cakṣur-dhātur yāvan na manovijñāna-dhātuḥ.}
&
\tib{མིག་གི་ཁམས་མེད་པ་ནས་ཡིད་ཀྱི་ཁམས་མེད། ཡིད་ཀྱི་རྣམ་པར་ཤེས་པའི་ཁམས་ཀྱི་བར་དུ་ཡང་མེད་དོ།}
\newline\rom{mig gi khams med pa nas yid kyi khams med/ yid kyi rnam par shes pa'i khams kyi bar du yang med do/}
\\
\hline
\engltrans{no eye-element, up to and including no mind-consciousness element.}
\hline
\hline

\segnum{7} 無無明，亦無無明盡，乃至無老死，亦無老死盡。
\newline\rom{Wú wúmíng, yì wú wúmíng jìn, nǎizhì wú lǎosǐ, yì wú lǎosǐ jìn.}
&
\deva{न-अविद्या न-अविद्या-क्षयो यावन् न जरा-मरणं न जरा-मरण-क्षयो।}
\newline\rom{Na-avidyā na-avidyā-kṣayo yāvan na jarā-maraṇaṃ na jarā-maraṇa-kṣayo.}
&
\tib{མ་རིག་པ་མེད། མ་རིག་པ་ཟད་པ་མེད་པ་ནས། རྒ་ཤི་མེད། རྒ་ཤི་ཟད་པའི་བར་དུ་ཡང་མེད་དོ།}
\newline\rom{ma rig pa med/ ma rig pa zad pa med pa nas/ rga shi med/ rga shi zad pa'i bar du yang med do/}
\\
\hline
\engltrans{No ignorance and no ending of ignorance, up to and including no old age and death, and no ending of old age and death.}
\hline
\hline

\segnum{8} 無苦集滅道，無智亦無得，以無所得故。
\newline\rom{Wú kǔ jí miè dào, wú zhì yì wú dé, yǐ wú suǒ dé gù.}
&
\deva{न दुःख-समुदय-निरोध-मार्गा। न ज्ञानं, न प्राप्तिर् न-अप्राप्तिः।}
\newline\rom{Na duḥkha-samudaya-nirodha-mārgā. Na jñānaṃ, na prāptir na-aprāptiḥ.}
&
\tib{སྡུག་བསྔལ་བ་དང་། ཀུན་འབྱུང་བ་དང་། འགོག་པ་དང་། ལམ་མེད། ཡེ་ཤེས་མེད། ཐོབ་པ་མེད། མ་ཐོབ་པ་ཡང་མེད་དོ།}
\newline\rom{sdug bsngal ba dang / kun 'byung ba dang / 'gog pa dang / lam med/ ye shes med/ thob pa med/ ma thob pa yang med do/}
\\
\hline
\engltrans{No suffering, origination, cessation, or path; no wisdom and no attainment, because there is nothing to attain.}
\hline
\hline

\segnum{9} 菩提薩埵，依般若波羅蜜多故。心無罣礙，無罣礙故。無有恐怖，遠離顚倒夢想，究竟涅槃。
\newline\rom{Pútísàduǒ, yī bōrě bōluómìduō gù. Xīn wú guà'ài, wú guà'ài gù. Wú yǒu kǒngbù, yuǎnlí diāndǎo mèngxiǎng, jiūjìng nièpán.}
&
\deva{तस्माच् छारिपुत्र अप्राप्तित्वाद् बोधिसत्त्वस्य प्रज्ञापारमितामाश्रित्य विहरत्यचित्तावरणः। चित्तावरण-नास्तित्वाद् अत्रस्त्रो विपर्यास-अतिक्रान्तो निष्ठा-निर्वाण-प्राप्तः।}
\newline\rom{Tasmāc Chāriputra aprāptitvād bodhisattvasya prajñāpāramitām āśritya viharaty acittāvaraṇaḥ. Cittāvaraṇa-nāstitvād atrastro viparyāsa-atikrānto niṣṭhā-nirvāṇa-prāptaḥ.}
&
\tib{ཤཱ་རིའི་བུ། དེ་ལྟ་བས་ན་བྱང་ཆུབ་སེམས་དཔའ་རྣམས་ཐོབ་པ་མེད་པའི་ཕྱིར། ཤེས་རབ་ཀྱི་ཕ་རོལ་ཏུ་ཕྱིན་པ་ལ་བརྟེན་ཅིང་གནས་ཏེ། སེམས་ལ་སྒྲིབ་པ་མེད་པས་འཇིགས་པ་མེད་དེ། ཕྱིན་ཅི་ལོག་ལས་ཤིན་ཏུ་འདས་ནས་མྱ་ངན་ལས་འདས་པའི་མཐར་ཕྱིན་ཏོ།}
\newline\rom{shA ri'i bu/ de lta bas na byang chub sems dpa' rnams thob pa med pa'i phyir/ shes rab kyi pha rol tu phyin pa la brten cing gnas te/ sems la sgrib pa med pas 'jigs pa med de/ phyin ci log las shin tu 'das nas mya ngan las 'das pa'i mthar phyin to/}
\\
\hline
\engltrans{Because bodhisattvas rely on prajñāpāramitā, their minds are without obstruction. Without obstruction, they have no fear, are far removed from inverted views and dreams, and attain ultimate nirvāṇa.}
\hline
\hline

\segnum{10} 三世諸佛，依般若波羅蜜多故，得阿耨多羅三藐三菩提。
\newline\rom{Sānshì zhū fó, yī bōrě bōluómìduō gù, dé ānòuduōluó sānmiǎo sānpútí.}
&
\deva{त्र्यध्व-व्यवस्थिताः सर्व-बुद्धाः प्रज्ञापारमितामाश्रित्यानुत्तरां सम्यक्-संबोधिमभिसंबुद्धाः।}
\newline\rom{Tryadhva-vyavasthitāḥ sarva-buddhāḥ prajñāpāramitām āśrityānuttarāṃ samyak-saṃbodhim abhisaṃbuddhāḥ.}
&
\tib{དུས་གསུམ་དུ་རྣམ་པར་བཞུགས་པའི་སངས་རྒྱས་ཐམས་ཅད་ཀྱང་ཤེས་རབ་ཀྱི་ཕ་རོལ་ཏུ་ཕྱིན་པ་ལ་བརྟེན་ནས་བླ་ན་མེད་པ་ཡང་དག་པར་རྫོགས་པའི་བྱང་ཆུབ་ཏུ་མངོན་པར་རྫོགས་པར་སངས་རྒྱས་སོ།}
\newline\rom{dus gsum du rnam par bzhugs pa'i sangs rgyas thams cad kyang shes rab kyi pha rol tu phyin pa la brten nas bla na med pa yang dag par rdzogs pa'i byang chub tu mngon par rdzogs par sangs rgyas so/}
\\
\hline
\engltrans{All buddhas of the three times, relying on prajñāpāramitā, attain anuttarā-samyak-saṃbodhi.}
\hline
\hline

\segnum{11} 故知般若波羅蜜多，是大神咒，是大明咒，是無上咒，是無等等咒，能除一切苦，眞實不虛故。
\newline\rom{Gùzhī bōrě bōluómìduō, shì dà shén zhòu, shì dà míng zhòu, shì wúshàng zhòu, shì wú děngděng zhòu, néng chú yīqiè kǔ, zhēnshí bù xū gù.}
&
\deva{तस्माज् ज्ञातव्यम्: प्रज्ञापारमिता महा-मन्त्रो महा-विद्या-मन्त्रो ऽनुत्तर-मन्त्रो ऽसमसम-मन्त्रः, सर्व-दुःख-प्रशमनः, सत्यम् अमिथ्यत्वात्।}
\newline\rom{Tasmāj jñātavyam: prajñāpāramitā mahā-mantro mahā-vidyā-mantro 'nuttara-mantro 'samasama-mantraḥ, sarva-duḥkha-praśamanaḥ, satyam amithyatvāt.}
&
\tib{དེ་ལྟ་བས་ན་ཤེས་རབ་ཀྱི་ཕ་རོལ་ཏུ་ཕྱིན་པའི་སྔགས། རིག་པ་ཆེན་པོའི་སྔགས། བླ་ན་མེད་པའི་སྔགས། མི་མཉམ་པ་དང་མཉམ་པའི་སྔགས། སྡུག་བསྔལ་ཐམས་ཅད་རབ་ཏུ་ཞི་བར་བྱེད་པའི་སྔགས། མི་རྫུན་པས་ན་བདེན་པར་ཤེས་པར་བྱ་སྟེ། ཤེས་རབ་ཀྱི་ཕ་རོལ་ཏུ་ཕྱིན་པའི་སྔགས་སྨྲས་པ།}
\newline\rom{de lta bas na shes rab kyi pha rol tu phyin pa'i sngags/ rig pa chen po'i sngags/ bla na med pa'i sngags/ mi mnyam pa dang mnyam pa'i sngags/ sdug bsngal thams cad rab tu zhi bar byed pa'i sngags/ mi rdzun pas na bden par shes par bya ste/ shes rab kyi pha rol tu phyin pa'i sngags smras pa/}
\\
\hline
\engltrans{Therefore know that prajñāpāramitā is the great mantra, the mantra of great illumination, the supreme mantra, the unequalled mantra, which can remove all suffering; it is true and not false.}
\hline
\hline

\segnum{12} 說般若波羅蜜多咒，卽說咒曰：揭帝揭帝，波羅揭帝，波羅僧揭帝，菩提僧莎訶。
\newline\rom{Shuō bōrě bōluómìduō zhòu, jí shuō zhòu yuē: Jiēdì jiēdì, bōluó jiēdì, bōluó sēng jiēdì, pútí sēng suōhē.}
&
\deva{प्रज्ञापारमितायामुक्तो मन्त्रः। तद्यथा: ॐ गते गते पारगते पारसंगते बोधि स्वाहा।}
\newline\rom{Prajñāpāramitāyām ukto mantraḥ. Tadyathā: Oṃ gate gate pāragate pārasaṃgate bodhi svāhā.}
&
\tib{ཏདྱ་ཐཱ། ཨོཾ་ག་ཏེ་ག་ཏེ་པཱ་ར་ག་ཏེ་པཱ་ར་སཾ་ག་ཏེ་བོ་དྷི་སྭཱ་ཧཱ།}
\newline\rom{tad+ya thA/ oM ga te ga te pA ra ga te pA ra saM ga te bo d+hi swA hA/}
\\
\hline
\engltrans{Therefore proclaim the prajñāpāramitā mantra, proclaim the mantra that says: Gate gate pāragate pārasaṃgate bodhi svāhā.}
\hline
\hline

\end{longtable}\section*{Critical Apparatus}\footnotesize\begin{multicols}{2}\var{1}~\wit{GRETIL}: Section I error (Attwood 2015): pañca-skandhās should be pañcaskandhāṃs (accusative). The skandhas are objects of vyavalokayati sma ('examined'), not nominative subjects. Also: vyavalokayati sma uses imperfect tense unusual for sūtra style.
\var{2}~\wit{Toh21}: T251 omits frame context; Tibetan preserves 'at that time' temporal marker
\var{3}~\wit{GRETIL}: The phrase 'yad rūpaṃ sā śūnyatā yā śūnyatā tad rūpam' has no Chinese counterpart and is absent from most Conze witnesses. Late interpolation not in Large Sutra parallels (Nattier 1992, Attwood 2017a). Original Aṣṭa reading was 'rūpam eva māyā māyaiva rūpam' (appearance is only illusion).
\var{4}~\wit{GRETIL}: Uses kṣaya (destruction) instead of standard nirodha (cessation) - evidence of back-translation (Nattier 1992)
\var{5}~\wit{Toh21}: Tibetan uses zad pa (exhaustion) matching Sanskrit kṣaya rather than standard nirodha - evidence for Nattier's back-translation thesis
\var{6}~\wit{GRETIL}: na-aprāptiḥ is a late interpolation not found in any Large Sutra version (Huifeng 2014). Should be 'na prāptir nābhisamayaḥ' (no attainment, no realization). The pair prāpti/abhisamaya refers to mārga/phala (path and fruition). Kumārajīva's T223 error was copied into the Heart Sutra (Attwood 2020a).
\var{7}~\wit{GRETIL}: Section VI error (Attwood 2018a, 2020a): Conze's wrongly placed full stop creates a sentence fragment. niṣṭhānirvāṇaḥ is a bahuvrīhi compound (masculine adjective) meaning 'one whose nirvāṇa is complete' - it describes the bodhisattva, not a separate neuter noun. Corrected reading removes the period: 'viharaty acittāvaraṇaś cittāvaraṇanāstitvād atrasto viparyāsātikrānto niṣṭhānirvāṇaḥ'.
\var{8}~\wit{GRETIL}: Epithets section (Attwood 2017b): Standard Prajñāpāramitā has three epithets (mahā-, anuttara-, asamasama-) representing vidyā, not mantra. Gilgit Pañc has 'mahāvidyeyaṃ...anuttareyaṃ vidyā...asamasameyaṃ vidyā'. Chinese míngzhòu 明咒 represents vidyā. The fourth epithet shénzhòu 神咒 may be a synonym for míngzhòu.
\var{9}~\wit{GRETIL}: Oṃ before 'gate gate' is absent from most Sanskrit manuscripts including the Hōryūji manuscript, and from all Chinese versions except T.257 (11th century Dānapāla translation). Late Tantric interpolation (Attwood 2024). Conze (1975) omits oṃ from his translation.
\var{10}~\wit{Fangshan}: The Fangshan reading 薩婆訶 (sàpóhē) is a more accurate transliteration of Sanskrit svāhā than T251's 僧莎訶 (sēng suōhē), which appears to be corrupt.
\end{multicols}
\end{document}
